\chapter{Validation}\label{ch:validation}

Do the procedures in this sequence: the checks on each run, the subsonic references,
the transonic and supersonic references, the flight, and the mesh convergence before you
publish a value.

\section{Checks on each run}\label{sec:run-checks}

\begin{longtable}{@{}P{0.24\linewidth} P{0.32\linewidth} P{0.36\linewidth}@{}}
\caption{Checks on each run.}\label{tab:run-checks}\\
\toprule
Check & Correct result & Cause if not correct \\
\midrule
\endfirsthead
\toprule
Check & Correct result & Cause if not correct \\
\midrule
\endhead
$C_L$, $C_m$ at $\alpha = 0$; $C_l$ (roll) & $\approx 0$ & The mesh is not symmetric, or the run did not converge. \\
\yplus{} on each wall patch & Near the estimate of \cref{sec:layers}, and equal on the meshes of a study & The first cell changed. \\
Minimum and maximum of $T$ & Physical: at most near the stagnation temperature & The start transient is not resolved. \\
Iterations & More than \ent{minIter}, and the criterion met & The run stopped before it was steady. \\
$C_N(\alpha)$ & Linear at small $\alpha$ & The run did not converge, or the fins are not resolved. \\
Machine record & No suspension; one CPU profile & The cost is not usable (\cref{sec:measurement}). \\
\bottomrule
\end{longtable}

\section{Subsonic flow: Barrowman}\label{sec:barrowman}

The slender-body equations of \citet{barrowman1967,barrowman1966}, which OpenRocket
uses, give the normal force and the centre of pressure of each component, referred to
the body cross-section and for each radian. \citet[ch.~2]{fleeman2012} gives equivalent
equations for the fins and the fin-body interference.

\begin{tabularx}{\linewidth}{@{}l L@{}}
\toprule
Component & $C_{N\alpha}$ \\
\midrule
Nose & 2, for any nose of revolution \\
Cylinder & $\approx 0$ \\
Boattail & $2\left[(d_\text{base}/d)^2 - 1\right]$ \\
Fins, $N$ fins & $K_{fb}\,\dfrac{4N(s/d)^2}{1 + \sqrt{1 + \left(\dfrac{2 l_m}{c_r + c_t}\right)^2}}$, \quad $K_{fb} = 1 + \dfrac{r}{s + r}$ \\
\bottomrule
\end{tabularx}

\medskip
$s$ is the exposed semi-span, $c_r$ and $c_t$ the root and tip chords, $l_m$ the length
of the mid-chord line, $r$ the body radius.

\begin{steps}
\item Run an angle-of-attack sweep at small angles.
\item Calculate the slope of $C_N = C_L\cos\alpha + C_D\sin\alpha$ against $\alpha$ in
  radians. $C_L$ is normal to the free stream, $C_N$ to the body axis.
\item Compare the slope with the sum of the components. For the breakdown, add one
  \ent{forceCoeffs} function object for each patch.
\end{steps}

\paragraph{Drag at zero angle.} Compare the components, not only the total. The
viscous part of \file{force.dat} against the skin friction $C_f S_\text{wet}/A_\text{ref}$
with $C_f = 0.0576\,Re_L^{-1/5}$ \citep[ch.~21]{schlichting1979} and the form factors
$1 + 1/(2L/d)$ (body) and $1 + 2t/c$ (fins); the base drag against
\citet[ch.~3]{hoerner1965}. $Re_L$ of each case is in \file{run.info}.

\section{Transonic and supersonic flow}\label{sec:supersonic-validation}

\begin{enumerate}
\item \textbf{Tutorial of the template.} Run the \cmd{rhoCentralFoam} tutorial
  \file{biconic25-55Run35} and compare it with its published data. This separates a
  configuration error from a mesh error.
\item \textbf{Oblique shock on a wedge.} The $\theta$--$\beta$--$M$ relation is exact
  \citep[ch.~4]{anderson2003}:
  \[
    \tan\theta = 2\cot\beta\,\frac{M^2\sin^2\beta - 1}{M^2(\gamma + \cos 2\beta) + 2}.
  \]
  The tutorial \file{wedge15Ma5} is this case.
\item \textbf{Flight.} \Cref{ch:flight}. The CFD must show the transonic drag rise
  of the flight. If it does not, the mesh is too coarse where the shocks occur.
\item \textbf{OpenRocket.} Semi-empirical: it finds gross errors, not small ones.
\end{enumerate}

\section{Mesh convergence}\label{sec:gci}

With $f_1$, $f_2$, $f_3$ the results on the fine, medium and coarse meshes, the
representative cell size $h = N^{-1/3}$ and the ratio $r = h_\text{coarse}/h_\text{fine}$,
the method of \citet{celik2008} gives the observed order $p$, the extrapolated value and
the grid convergence index:
\[
  p = \frac{\ln\left[(f_3 - f_2)/(f_2 - f_1)\right]}{\ln r},\qquad
  f_\text{ext} = f_1 + \frac{f_1 - f_2}{r^p - 1},\qquad
  \text{GCI}_{12} = 1.25\,\frac{|(f_1 - f_2)/f_1|}{r^p - 1}.
\]
With unequal ratios, use the iterative form of \citet{celik2008}. If
$(f_3 - f_2)/(f_2 - f_1) < 0$, the convergence is oscillatory and the equations do not
apply \citep{roache1994}.

\section{Record}

For each published value keep: \file{meshDict}, \file{meshInfo} and
\file{log.checkMesh} of the mesh; the STL file or its commit; \file{run.info},
\file{flowConditions} and \file{config.json} of the run; and \file{postProcessing/}.
With these files another person can make the result again.
